\documentclass{article}
\usepackage[a4paper,margin=2cm]{geometry}
\usepackage{amsmath, amssymb, amsthm}
\usepackage{booktabs}
\newtheorem{proposition}{Proposition}
\newtheorem{lemma}[proposition]{Lemma}
\newtheorem{corollary}[proposition]{Corollary}
\newtheorem{conjecture}[proposition]{Conjecture}
\theoremstyle{definition}
\newtheorem*{definition*}{Definition}
\theoremstyle{remark}
\newtheorem*{remark}{Remark}

\title{A condition for a local maximum of $n\mapsto w_{n,p}$}
\author{}
\date{2026-04-17}

\begin{document}
\maketitle

\section*{Setting}

Consider the all-heads-wins game with parameter $p\in(0,1)$: in each round,
all remaining coins are tossed, at least one coin must be set aside, and
the game is won iff every set-aside coin shows heads. Let $w_{n,p}$ denote
the optimal winning probability starting with $n$ coins, $w_{0,p}:=1$.
By elementary conditioning on the number $j$ of tails in the first round,
\begin{equation}\label{eq:bellman}
  w_{n,p}\;=\;p^{n}+\sum_{j=1}^{n-1}\binom{n}{j}p^{\,n-j}(1-p)^{j}\,
  \max_{j\le m\le n-1}w_{m,p},\qquad n\ge 1.
\end{equation}
Throughout, write
\[
  M_{n,j}\;:=\;\max_{j\le m\le n-1}w_{m,p}\qquad(1\le j\le n-1).
\]

\section*{Local maximum}

\begin{definition*}
We say that $n\mapsto w_{n,p}$ has a \emph{strict local maximum} at
$n^{\star}\ge 2$ if
\[
  w_{n^{\star}-1,p}\;<\;w_{n^{\star},p}\;>\;w_{n^{\star}+1,p}.
\]
\end{definition*}

\subsection*{Definitional inequalities}

Directly from the definition, $n^{\star}$ is a strict local maximum iff
\begin{equation}\label{eq:def}
  w_{n^{\star},p}-w_{n^{\star}-1,p}>0
  \quad\text{and}\quad
  w_{n^{\star},p}-w_{n^{\star}+1,p}>0.
\end{equation}

\subsection*{Reduction to strategy (B)}

Let $b_{n,p}$ denote the winning probability under strategy~(B) (set aside
every head). It satisfies the same recursion \eqref{eq:bellman} with the
$\max_{j\le m\le n-1}w_{m,p}$ replaced by $b_{j,p}$:
\begin{equation}\label{eq:brec}
  b_{n,p}\;=\;p^{n}+\sum_{j=1}^{n-1}\binom{n}{j}p^{\,n-j}(1-p)^{j}\,b_{j,p},
  \qquad b_{0,p}:=1.
\end{equation}

\begin{lemma}[Reduction]\label{lem:reduction}
If $(w_{m,p})_{m=1}^{N}$ is monotone non-increasing, then
$w_{m,p}=b_{m,p}$ for every $m\in\{1,\dots,N\}$.
\end{lemma}

\begin{proof}
Induction on $m$. For $m=1$ both sides equal $p$. For $2\le m\le N$, the
monotonicity hypothesis gives
$\max_{j\le m'\le m-1}w_{m',p}=w_{j,p}=b_{j,p}$ by the inductive
hypothesis, so \eqref{eq:bellman} at index $m$ reduces exactly to
\eqref{eq:brec} at index $m$. \qed
\end{proof}

\begin{corollary}[First non-monotonicity of $w$ via $b$]\label{cor:firstmax}
Let $n^{\star}\ge 2$ be the smallest integer with
$b_{n^{\star},p}>b_{n^{\star}-1,p}$. Then $w_{m,p}=b_{m,p}$ for every
$m\le n^{\star}$, and
\begin{equation}\label{eq:first-uptick}
  w_{n^{\star},p}\;>\;w_{n^{\star}-1,p}.
\end{equation}
In particular $n\mapsto w_{n,p}$ is not monotone non-increasing.
\end{corollary}

\begin{proof}
By the minimality of $n^{\star}$ in the statement (and by the
equalities $w_{m,p}=b_{m,p}$ proved below for $m<n^{\star}$),
$(b_{m,p})_{m=1}^{n^{\star}-1}$ is monotone non-increasing. If
$(w_{m,p})_{m=1}^{n^{\star}-1}$ is also monotone non-increasing,
Lemma~\ref{lem:reduction} gives $w_{m,p}=b_{m,p}$ for $m<n^{\star}$; we
verify this by induction. For $m=1$ both equal $p$. Assume
$w_{m',p}=b_{m',p}$ for $m'<m\le n^{\star}-1$; then
$b_{1,p}\ge\dots\ge b_{m-1,p}$ translates to
$w_{1,p}\ge\dots\ge w_{m-1,p}$, so
$M_{m,j}=w_{j,p}=b_{j,p}$ in \eqref{eq:bellman}, and~\eqref{eq:bellman}
reduces to~\eqref{eq:brec}, giving $w_{m,p}=b_{m,p}$. The same
argument at $m=n^{\star}$ gives $w_{n^{\star},p}=b_{n^{\star},p}$. In
particular $w_{n^{\star},p}=b_{n^{\star},p}>b_{n^{\star}-1,p}=w_{n^{\star}-1,p}$.
\qed
\end{proof}

\begin{remark}
Corollary~\ref{cor:firstmax} does \emph{not} claim that $n^{\star}$ is
a strict local maximum of $n\mapsto w_{n,p}$. Indeed, at $p=0.45$ the
first $b$-uptick occurs at $n^{\star}=7$ (with $b_{7}\approx 0.4111$),
while the first strict local maximum of $w$ occurs only at $n=15$;
between $n^{\star}$ and $15$ the sequence $w_{n,p}$ continues to
increase (through further $b$-upticks). Pinning down the location of
a local maximum requires running the Bellman recursion past
$n^{\star}$.
\end{remark}

\section*{Uniform upper bounds on $w_{n,p}$ for $p<\tfrac12$}

The Bellman value never exceeds $w_{1,p}=p$ for $n\ge 2$, and never
exceeds $w_{2,p}$ for $n\ge 3$. Both bounds follow by direct
induction on the Bellman.

\begin{proposition}[$w_{n,p}<p$ for every $n\ge 2$, $p<\tfrac12$]\label{prop:wn-bounded-by-p}
Let $p\in(0,\tfrac12)$ and $q:=1-p$. Then for every $n\ge 2$,
\begin{equation}\label{eq:wn-bound}
  w_{n,p}\;\le\;p-pq\bigl(q^{n-1}-p^{n-1}\bigr)\;<\;p\;=\;w_{1,p},
\end{equation}
with equality at $n=2$.
\end{proposition}

\begin{proof}
By induction on $n$.

\emph{Base $n=2$.} Strategies (A) and (B) coincide on $n=2$ and give
$w_{2,p}=p^2(3-2p)$. Hence
$p-w_{2,p}=p\bigl(1-3p+2p^{2}\bigr)=p(1-p)(1-2p)=pq(q-p)=pq(q^{1}-p^{1})$.
This is positive for $p\in(0,\tfrac12)$ and matches~\eqref{eq:wn-bound}
with equality.

\emph{Induction step} ($n\ge 3$). Assume $w_{m,p}<p$ for every
$m\in\{2,\dots,n-1\}$, so that
$\max_{j\le m\le n-1}w_{m,p}=w_{1,p}=p$ when $j=1$ (the unique
maximizer is $m=1$) and $\max_{j\le m\le n-1}w_{m,p}<p$ when $j\ge 2$
(the range excludes $1$). Substituting in the Bellman~\eqref{eq:bellman}
and using strictness for $j\ge 2$ (the sum is non-empty since
$n\ge 3$):
\begin{align*}
  w_{n,p}
  &\;<\;p^n+\binom{n}{1}p^{n-1}q\cdot p
       +p\sum_{j=2}^{n-1}\binom{n}{j}p^{n-j}q^{j}\\
  &\;=\;p^n+np^nq+p\bigl(1-p^n-np^{n-1}q-q^n\bigr)\\
  &\;=\;p+p^nq-pq^n
   \;=\;p-pq\bigl(q^{n-1}-p^{n-1}\bigr),
\end{align*}
where the second equality uses
$\sum_{j=2}^{n-1}\binom{n}{j}p^{n-j}q^{j}=1-p^n-np^{n-1}q-q^n$ (from
the binomial theorem). The right-hand side is $<p$ because $q>p$ for
$p<\tfrac12$, hence $q^{n-1}>p^{n-1}$.\qed
\end{proof}

The same induction template, with $w_{2,p}$ in place of $w_{1,p}$,
yields a strictly tighter bound.

\begin{proposition}[$w_{n,p}<w_{2,p}$ for every $n\ge 3$, $p<\tfrac12$]\label{prop:wn-bounded-by-w2}
Let $p\in(0,\tfrac12)$. For every $n\ge 3$,
\begin{equation}\label{eq:wn-bound2}
  w_{n,p}\;\le\;w_{2,p}-p^{2}q^{2}\,G_{n}(p)\;<\;w_{2,p},
\end{equation}
where
\begin{equation}\label{eq:Gn-def}
  G_{n}(p)\;:=\;q^{\,n-2}(3-2p)-p^{\,n-2}\bigl[(n+1)-2p(n-1)\bigr].
\end{equation}
\end{proposition}

\begin{proof}
\emph{Step~1: setup.} By induction on $n$ assume $w_{m,p}\le w_{2,p}$
for $3\le m\le n-1$. Combined with
Proposition~\ref{prop:wn-bounded-by-p} ($w_{m,p}<p=w_{1,p}$ for
$m\ge 2$) the Bellman maxes split:
\begin{itemize}
  \item $j=1$: $\max_{1\le m\le n-1}w_{m,p}=w_{1,p}=p$;
  \item $j=2$: $\max_{2\le m\le n-1}w_{m,p}=w_{2,p}$ (by the induction
        hypothesis, $w_{m,p}\le w_{2,p}$ for $m\ge 3$);
  \item $j\ge 3$: $\max_{j\le m\le n-1}w_{m,p}\le w_{2,p}$ (induction
        hypothesis, strict for $n\ge 4$).
\end{itemize}
Substituting in~\eqref{eq:bellman} and using
$\sum_{j=2}^{n-1}\binom{n}{j}p^{n-j}q^{j}=1-p^n-np^{n-1}q-q^n$:
\[
  w_{n,p}\;\le\;\bar w_n\;:=\;p^n(1+nq)+w_{2,p}\bigl(1-p^n-np^{n-1}q-q^n\bigr).
\]

\emph{Step~2: reduction to $G_{n}(p)>0$.} Using
$p-w_{2,p}=pq(1-2p)$ and grouping,
$\bar w_n-w_{2,p}=-p^{2}q^{2}\,G_{n}(p)$, so $\bar w_n<w_{2,p}$
iff $G_n(p)>0$.

\emph{Step~3: $G_{n}(p)>0$ on $(0,\tfrac12)$ for $n\ge 3$.} By a
second induction on $n$, with the algebraic identity
\begin{equation}\label{eq:Gn-rec}
  G_{n}(p)-q\,G_{n-1}(p)\;=\;p^{\,n-3}(1-2p)\bigl[3p+n(q-p)\bigr]
  \qquad(n\ge 4),
\end{equation}
which is verified directly. For $p\in(0,\tfrac12)$, both factors
$(1-2p)$ and $[3p+n(q-p)]$ are positive. Base case:
$G_{3}(p)=q\bigl[(3-2p)-4p\bigr]=3q(1-2p)>0$. Inductive step:
$G_{n}(p)=qG_{n-1}(p)+p^{n-3}(1-2p)[3p+n(q-p)]>qG_{n-1}(p)>0$.

\emph{Step~4: conclude.} $w_{n,p}\le\bar w_{n}=w_{2,p}-p^{2}q^{2}G_{n}(p)<w_{2,p}$.\qed
\end{proof}

\begin{remark}
The cancellation in Step~2 — which collapses the
$\binom{n}{2}p^{n-2}q^{2}$ terms — is what allows the induction to
bootstrap from $w_{2,p}$ rather than just from $w_{1,p}=p$. The
recursion~\eqref{eq:Gn-rec} is the algebraic engine: it says
$G_{n}-qG_{n-1}$ equals an explicitly positive quantity on
$(0,\tfrac12)$, so the strict positivity of $G_{n}$ propagates upward
in $n$.
\end{remark}

Empirically, for every tested $p\in(0.05,0.5)$ at finite horizon
$n_{\max}\le 300$, the Bellman sequence has many strict local maxima,
all with values $w_{N,p}<w_{2,p}$ — consistent with
Propositions~\ref{prop:wn-bounded-by-p}--\ref{prop:wn-bounded-by-w2}.
The set of running-maximum indices of $(w_{n,p})_{n\ge 1}$ is
exactly $\{1\}$ for every tested $p\in(0,\tfrac12)$.

\subsection*{Attempt to extend to $w_{n,p}<w_{3,p}$}

The natural next step is the bound $w_{n,p}<w_{3,p}$ for $n\ge 4$,
$p<\tfrac12$. Numerically the bound holds for every tested
$(n,p)$. The same induction template — but now with a fourth bucket
$j=3$ — gives, after grouping,
\begin{equation}\label{eq:bar-w-w3}
  \bar w_n - w_{3,p}\;=\;p^{3}q^{2}\,(2p-1)\,R_n(p),
\end{equation}
where $R_n$ is an explicit polynomial in $p$ of degree $n$ with
$R_n(0)=13$ and $R_n(\tfrac12)\ge 0$. The bound $\bar w_n\le w_{3,p}$
is therefore equivalent (since $(2p-1)<0$ on $(0,\tfrac12)$) to
$R_n(p)>0$ on $(0,\tfrac12)$.

\begin{proposition}[partial result on $w_n<w_3$]\label{prop:wn-bounded-by-w3-partial}
For every $p\in(0,\tfrac12)$ and every $n\in\{4,\dots,12\}$,
$w_{n,p}<w_{3,p}$.
\end{proposition}

\emph{Sketch.} Same as Proposition~\ref{prop:wn-bounded-by-w2}: the
Bellman maxes split as $j=1\Rightarrow w_1$, $j=2\Rightarrow w_2$,
$j=3\Rightarrow w_3$, $j\ge 4\Rightarrow$ value bounded by $w_3$ via
the induction hypothesis. Substituting and identifying the
sum~$\sum_{j=3}^{n-1}\binom{n}{j}p^{n-j}q^{j}=1-p^n-q^n-np^{n-1}q-\binom{n}{2}p^{n-2}q^{2}$
yields~\eqref{eq:bar-w-w3}. The polynomial $R_n$ is then verified to
be positive on the entire interval $(0,\tfrac12)$ by direct
real-roots computation for each $n\in\{4,\dots,12\}$.\qed

\begin{remark}[Breakdown for $n\ge 13$ near $p=\tfrac12$]
Unlike $G_n$ in Proposition~\ref{prop:wn-bounded-by-w2}, the
polynomial $R_n$ does \emph{not} satisfy a clean positivity-preserving
recursion of the form $R_n - qR_{n-1}=(\text{positive})$. Numerically,
for $n\ge 13$ the polynomial $R_n$ acquires a single real root
$p_n\in(0,\tfrac12)$ that drifts towards $\tfrac12$ from below as $n$
grows ($p_{13}\approx 0.4978$, $p_{20}\approx 0.4911$,
$p_{24}\approx 0.4900$). Hence the bound $\bar w_n\le w_{3,p}$
\emph{fails} on the subinterval $(p_n,\tfrac12)$ for $n\ge 13$,
even though numerically the sharper claim $w_{n,p}<w_{3,p}$ continues
to hold throughout. Closing the gap for $n\ge 13$ requires a tighter
induction (e.g.\ chaining further bounds $w_{m,p}\le w_{4,p},w_{5,p},\dots$)
to refine the $j\ge 4$ contributions.
\end{remark}

\section*{A tool for finding local maxima}

Corollary~\ref{cor:firstmax} reduces the problem of locating a local
maximum of the \emph{Bellman} value function $w$ to the (much simpler)
problem of detecting non-monotonicity of the \emph{strategy-(B)} sequence
$b$. Strategy (B) has no $\max$ in its recursion~\eqref{eq:brec}; it is
linear in the previous values. We record a few identities and structural
observations that make the tool tangible.

Throughout, set $q:=1-p$.

\subsection*{A probabilistic identity}

\begin{proposition}[Binomial-mixture identity]\label{prop:binmix}
For every $n\ge 1$ and $p\in(0,1)$,
\begin{equation}\label{eq:binmix}
  (1+q^{n})\,b_{n,p}\;=\;\mathbb E\bigl[b_{J,p}\bigr],
  \qquad J\sim\mathrm{Binomial}(n,q).
\end{equation}
Equivalently, $\mathbb E\bigl[b_{J,p}-b_{n,p}\bigr]=q^{n}b_{n,p}$.
\end{proposition}

\begin{proof}
Starting from the recursion \eqref{eq:brec} with $b_{0,p}=1$ and noting
that the $j=n$ term in the binomial expansion equals $q^{n}b_{n,p}$,
\begin{align*}
  b_{n,p}
  &=\sum_{j=0}^{n-1}\binom{n}{j}p^{\,n-j}q^{j}\,b_{j,p}
  \;=\;\sum_{j=0}^{n}\binom{n}{j}p^{\,n-j}q^{j}\,b_{j,p}-q^{n}b_{n,p}\\
  &=\mathbb E\bigl[b_{J,p}\bigr]-q^{n}b_{n,p}.\qed
\end{align*}
\renewcommand{\qed}{}
\end{proof}

\subsection*{Polynomial structure}

Writing $b_{n,p}=p^{n}P_{n}(q)$, the recursion~\eqref{eq:brec} becomes
\begin{equation}\label{eq:Prec}
  P_{n}(q)\;=\;\sum_{j=0}^{n-1}\binom{n}{j}q^{j}\,P_{j}(q),\qquad P_{0}:=1.
\end{equation}

\begin{proposition}\label{prop:Pstruct}
$P_{n}(q)$ is a polynomial in $q$ of degree $\binom{n}{2}$ with
\emph{non-negative integer} coefficients. Moreover
\[
  P_{n}(0)=1,\qquad P_{n}(1)=F_{n},\qquad P_{n}(\tfrac12)=2^{\,n-1},
\]
where $F_{n}=1,1,3,13,75,541,\dots$ is the $n$-th Fubini number
(ordered Bell number, \emph{OEIS~A000670}).
\end{proposition}

\begin{proof}
The recursion~\eqref{eq:Prec} is linear with non-negative integer
coefficients and $P_{0}=1$, so induction gives
non-negative integer coefficients. The degree recursion
$\deg P_{n}=(n-1)+\deg P_{n-1}$, with $\deg P_{1}=0$, gives
$\deg P_{n}=\binom{n}{2}$.

$P_{n}(0)$: only the $j=0$ term contributes at $q=0$, giving
$\binom{n}{0}P_{0}(0)=1$.

$P_{n}(1)$: at $q=1$, \eqref{eq:Prec} becomes
$P_{n}(1)=\sum_{j=0}^{n-1}\binom{n}{j}P_{j}(1)$, which is the defining
recursion of the Fubini numbers (with $F_{0}=1$).

$P_{n}(1/2)$: $b_{n,1/2}=\tfrac12$ for all $n\ge 1$, so
$P_{n}(1/2)=b_{n,1/2}/(1/2)^{n}=2^{\,n-1}$. \qed
\end{proof}

\subsection*{An exponential generating function}

\begin{proposition}
Let $G(x):=\sum_{n\ge 0}b_{n,p}\,x^{n}/n!$. Then
\begin{equation}\label{eq:egf}
  G(x)\;=\;1+(e^{px}-1)\,G(qx)
  \;=\;\sum_{k\ge 0}\prod_{i=0}^{k-1}\bigl(e^{pq^{i}x}-1\bigr).
\end{equation}
\end{proposition}

\begin{proof}
Multiply the recursion~\eqref{eq:brec} by $x^{n}/n!$ and sum over
$n\ge 1$; exchange orders of summation and identify
$\sum_{n>j}(px)^{n-j}/(n-j)!=e^{px}-1$. Iterating $G(x)=1+(e^{px}-1)G(qx)$
and noting that the remainder term
$\prod_{i=0}^{k-1}(e^{pq^{i}x}-1)\,G(q^{k}x)\to 0$ as $k\to\infty$
(since $e^{pq^{i}x}-1\sim pq^{i}x$ for large $i$), the product expansion
follows. \qed
\end{proof}

\subsection*{Small-$p$ heuristic}

At leading order in $p\to 0^{+}$, \eqref{eq:Prec} evaluated at $q\to 1$
gives $P_{n}(1)=F_{n}$ and hence
\[
  b_{n,p}\;\sim\;F_{n}\,p^{n}\qquad(p\to 0^{+},\ n\text{ fixed}).
\]
Fubini numbers satisfy $F_{n}/F_{n-1}\to n/\log 2$, so
\[
  \frac{b_{n,p}}{b_{n-1,p}}\;\sim\;p\cdot\frac{F_{n}}{F_{n-1}}\;\sim\;\frac{np}{\log 2}
  \qquad(p\to 0^{+},\ n\text{ fixed}).
\]
This formally predicts $b_{n,p}>b_{n-1,p}$ for $n>(\log 2)/p$. The
prediction is \emph{qualitatively} right but quantitatively much too
optimistic: the Taylor expansion in $p$ is pointwise (fixed $n$, $p\to 0$),
whereas the first monotonicity failure $n_{\mathrm{fail}}(p)$ diverges
with $n$ and the corrections of order $p^{n+1},p^{n+2},\dots$ cannot be
ignored. Empirically $n_{\mathrm{fail}}(p)\asymp p^{-2.2}$, not $p^{-1}$.

\section*{Toward a general proof}

\subsection*{Reformulation}

Combining the tool of Corollary~\ref{cor:firstmax} with
Proposition~\ref{prop:Pstruct}, our open conjecture admits the following
purely analytic reformulation.

\begin{conjecture}\label{conj:main}
For every $p\in(0,\tfrac12)$, the sequence $(b_{n,p})_{n\ge 1}$ is
\emph{not} monotone non-increasing in $n$. Equivalently, for every
$q\in(\tfrac12,1)$ there exists $n\ge 2$ such that
\begin{equation}\label{eq:conj}
  (1-q)\,P_{n}(q)\;>\;P_{n-1}(q),
\end{equation}
where $P_{n}$ is the polynomial of Proposition~\ref{prop:Pstruct}.
\end{conjecture}

By Corollary~\ref{cor:firstmax}, Conjecture~\ref{conj:main} implies that
$n\mapsto w_{n,p}$ has a strict local maximum for every $p\in(0,\tfrac12)$.

\subsection*{What is known}

We collect what can be proved toward Conjecture~\ref{conj:main}.

\emph{Boundary.} At $q=\tfrac12$ (i.e.\ $p=\tfrac12$),
Proposition~\ref{prop:Pstruct} gives
$(1-q)P_{n}(q)=\tfrac12\cdot 2^{n-1}=2^{n-2}=P_{n-1}(\tfrac12)$ for every
$n\ge 2$, so~\eqref{eq:conj} is an equality — the symmetric case is the
boundary at which (B) is marginally monotone ($b_{n,1/2}=\tfrac12$ for
all $n$).

\emph{At $q=1$.} $P_{n}(1)=F_{n}$, Fubini numbers, and $F_{n}/F_{n-1}\to\infty$,
so $(1-q)P_{n}(q)/P_{n-1}(q)\to\infty$ as $q\uparrow 1$ for every
$n\ge 2$; the inequality \eqref{eq:conj} is (trivially) satisfied in a
suitable limit.

\emph{Known specific $p$.} For $p=21/50$, the exact rational calculation
of \cite{proof_p042} gives $w_{9}>w_{8}$; by Corollary~\ref{cor:firstmax}
this is equivalent to $b_{9,21/50}>b_{8,21/50}$, i.e.~\eqref{eq:conj}
holds at $q=29/50$, $n=9$.

\emph{Asymptotic limit of $b$.} If $(b_{n,p})$ were monotone
non-increasing for all $n\ge 1$, then $b_{n,p}\downarrow L$ for some
$L\ge 0$, and Proposition~\ref{prop:binmix} together with bounded
convergence gives $(1+q^{n})b_{n,p}=\mathbb E[b_{J,p}]\to L$; the limit
relation is therefore self-consistent (no direct contradiction).

\subsection*{A reduced identity for the forward difference}

Let $\Delta b_{n}:=b_{n,p}-b_{n-1,p}$. From \eqref{eq:binmix} at indices
$n$ and $n-1$,
\begin{align*}
  (1+q^{n})b_{n,p}-(1+q^{n-1})b_{n-1,p}
  &=\mathbb E\bigl[b_{J_{n},p}\bigr]-\mathbb E\bigl[b_{J_{n-1},p}\bigr]\\
  &=q\,\mathbb E\bigl[b_{J_{n-1}+1,p}-b_{J_{n-1},p}\bigr]\\
  &=q\,\mathbb E\bigl[\Delta b_{J_{n-1}+1}\bigr],
\end{align*}
using the coupling $J_{n}\stackrel{d}{=}J_{n-1}+\mathrm{Bern}(q)$. The
left-hand side expands as $\Delta b_{n}+q^{n-1}(qb_{n,p}-b_{n-1,p})$, so
\begin{equation}\label{eq:forwarddiff}
  \boxed{\;\Delta b_{n}
  \;=\;q\,\mathbb E\bigl[\Delta b_{J_{n-1}+1}\bigr]
  +q^{n-1}\bigl(b_{n-1,p}-q\,b_{n,p}\bigr).\;}
\end{equation}
Identity~\eqref{eq:forwarddiff} expresses the target quantity
$\Delta b_{n}$ as a \emph{binomial-weighted average} of earlier
$\Delta b_{\bullet}$'s plus an explicit positive correction of order
$q^{n-1}$. If $(b_{n,p})$ were monotone non-increasing, then
$\Delta b_{\bullet}\le 0$, the first term of~\eqref{eq:forwarddiff}
would be $\le 0$, and the second would be
$q^{n-1}(b_{n-1,p}-qb_{n,p})\ge q^{n-1}(1-q)b_{n-1,p}>0$.

Thus Conjecture~\ref{conj:main} reduces to the following.

\begin{conjecture}[equivalent to Conjecture~\ref{conj:main}]\label{conj:equiv}
For every $p\in(0,\tfrac12)$ there exists $n$ such that
\[
  q\,\mathbb E\bigl[(-\Delta b_{J_{n-1}+1})\bigr]
  \;<\;q^{n-1}\bigl(b_{n-1,p}-q\,b_{n,p}\bigr).
\]
\end{conjecture}

Both sides of the inequality tend to $0$, but the right-hand side decays
like $q^{n-1}\cdot p\cdot b_{n-1,p}$ (still assuming monotone
non-increasing, so $b_{n-1,p}\ge L\ge 0$), whereas the left-hand side is
a binomial average of past decrements; a precise rate comparison is the
missing step.

\subsection*{Status}

Conjecture~\ref{conj:main} — and hence the analytic existence of a
strict local maximum of $n\mapsto w_{n,p}$ for every $p\in(0,\tfrac12)$
— remains open. The attack via \eqref{eq:forwarddiff} is the cleanest
route so far: it requires only control on the \emph{sizes} of past
decrements $|\Delta b_{j}|$ weighted by a binomial measure, compared
against an explicit exponentially small positive term. A quantitative
estimate of the form
\[
  \sum_{j=1}^{n-1}\binom{n-1}{j}p^{\,n-1-j}q^{j}\cdot|\Delta b_{j}|
  \;=\;o\bigl(q^{n-1}\cdot p\cdot b_{n-1,p}\bigr)
\]
along some subsequence $n\to\infty$ would close the argument.

\begin{thebibliography}{9}
\bibitem{proof_p042}\emph{An explicit $(n,k)$ at which the optimal strategy is neither (A) nor (B)},
  \texttt{Manuscript/proof\_intermediate\_p042.tex}, 2026.
\end{thebibliography}

\end{document}
